% Ausblick Einheit 4 von 4 · Satz vom Nullprodukt (Katalog Einheit 2; OS Kl. 10)
\hypertarget{e4}{}\einheitenkopf{Ausblick Einheit 4 von 4 · Satz vom Nullprodukt}
\zweigzeile{Hier lernst du, Gleichungen, bei denen ein Produkt null ist, ohne Formel zu lösen · kommt nächstes Jahr (am Gymnasium je nach Buch schon seit Klasse 8) · P10 · baut auf: Normalform und p-q-Formel (Einheit 2), Ausmultiplizieren und Ausklammern, Einsetzen}
\renewcommand{\theaufgabe}{\arabic{aufgabe}}\setcounter{aufgabe}{27}

\begin{aufgabe}{Ich erkenne, ob bei einem Produkt rechts null steht. Kreuze an. Du musst nichts lösen.}
\begin{geruest}
\gz{a}{$x \cdot (x + 4) = 0$}{\kreuz{rechts steht null}\kreuz{rechts steht eine andere Zahl}}
\gz{b}{$(x - 2) \cdot (x + 5) = 12$}{\kreuz{rechts steht null}\kreuz{rechts steht eine andere Zahl}}
\gz{c}{$(x + 1) \cdot (x - 6) = 0$}{\kreuz{rechts steht null}\kreuz{rechts steht eine andere Zahl}}
\gz{d}{$x \cdot (x - 3) = 10$}{\kreuz{rechts steht null}\kreuz{rechts steht eine andere Zahl}}
\gz{e}{$3x \cdot (x + 2) = 0$}{\kreuz{rechts steht null}\kreuz{rechts steht eine andere Zahl}}
\end{geruest}
\end{aufgabe}

\begin{aufgabe}{Ich kann eine Gleichung lösen, bei der ein Produkt null ist. Ein Produkt ist null, wenn ein Faktor null ist. Setze jeden Faktor null und löse.}
\beispiel{(x - 2) \cdot (x - 5) &= 0 \\ x - 2 &= 0 \quad\text{oder}\quad x - 5 = 0 \\ x_1 &= 2, \quad x_2 = 5}
\begin{gleichungsraster}[2]
\gl{(x - 1) \cdot (x - 4) = 0} & \gl{(x - 3) \cdot (x - 6) = 0} \\ \rule{0pt}{2mm} & \\
\gl{(x - 2) \cdot (x - 7) = 0} & \gl{(x - 5) \cdot (x - 8) = 0} \\ \rule{0pt}{2mm} & \\
\gl{(x + 3) \cdot (x - 1) = 0} & \gl{x \cdot (x + 6) = 0} \\ \rule{0pt}{2mm} & \\
\gl{(x + 2) \cdot (x + 2) = 0} & \\ \rule{0pt}{2mm} & \\
\end{gleichungsraster}
Prüfe durch Einsetzen, ob die Zahl eine Lösung ist. Rechne die Klammern zuerst.
\begin{geruest}
\gz{h}{$x = -3$ \ in \ $(x + 5) \cdot (x - 2) = -10$}{\kreuz{(wA)}\kreuz{(fA)}}
\end{geruest}
\end{aufgabe}

\begin{aufgabe}{Ich kann $x$ ausklammern und dann jeden Faktor null setzen. Klammere $x$ aus und löse.}
\beispiel{x^2 - 5x &= 0 \\ x \cdot (x - 5) &= 0 \\ x &= 0 \quad\text{oder}\quad x - 5 = 0 \\ x_1 &= 0, \quad x_2 = 5}
\begin{gleichungsraster}[3]
\gl{x^2 + 3x = 0} & \gl{x^2 - 6x = 0} \\ \rule{0pt}{2mm} & \\
\gl{2x^2 + 10x = 0} & \\ \rule{0pt}{2mm} & \\
\end{gleichungsraster}
Hier steht rechts keine Null. Multipliziere aus, bringe die Gleichung in Normalform und löse sie mit der p-q-Formel.
\begin{gleichungsraster}[4]
\gl{x \cdot (x + 3) = 28} & \\ \rule{0pt}{2mm} & \\
\end{gleichungsraster}
\begin{teile}
\teil Zeige durch Ausmultiplizieren, dass $x^2 - 9x + 20 = (x - 4) \cdot (x - 5)$ ist. Löse dann $x^2 - 9x + 20 = 0$ ohne Formel.
\teil Kreuze an, welche Zahl die Gleichung $x \cdot (x + 7) = -12$ erfüllt. (P10 2025)\\ \kreuz{$3$}\kreuz{$4$}\kreuz{$-3$}\kreuz{$-12$}
\end{teile}
\end{aufgabe}

\begin{aufgabe}{Ich finde den Fehler beim Lösen einer Produktgleichung.}
\begin{teile}
\teil Lukas löst $x^2 = 4x$ so:
\rechnung{x^2 &= 4x &&\mid :x \\ x &= 4}
Finde den Fehler, benenne ihn und korrigiere die Rechnung.
\end{teile}
Löse selbst.
\begin{gleichungsraster}[3]
\gl{x^2 = 9x} & \\ \rule{0pt}{2mm} & \\
\end{gleichungsraster}
\begin{teile}
\teil Mara löst $x \cdot (x - 4) = 12$ so:
\rechnung{x \cdot (x - 4) &= 12 \\ x &= 12 \quad\text{oder}\quad x - 4 = 12 \\ x_1 &= 12, \quad x_2 = 16}
Finde den Fehler, benenne ihn und korrigiere die Rechnung.
\end{teile}
Löse selbst.
\begin{gleichungsraster}[4]
\gl{x \cdot (x - 3) = 18} & \\ \rule{0pt}{2mm} & \\
\end{gleichungsraster}
\end{aufgabe}

\begin{aufgabe}{Ich kann begründen, warum ein Produkt null ist und warum man nicht durch $x$ teilen darf.}
\begin{teile}
\teil Begründe ohne Ausmultiplizieren, warum $x = 9$ die Gleichung $(x - 9) \cdot (x + 4) = 0$ löst.
\teil Begründe, warum man bei $x^2 = 11x$ nicht durch $x$ teilen darf.
\end{teile}
\end{aufgabe}

\begin{aufgabe}{Ich kann mit dem Nullprodukt eine Frage zu einem Wurf beantworten. Ein Ball wird vom Boden aus senkrecht nach oben geworfen. Nach $t$ Sekunden ist er $15t - 5t^2$ Meter hoch.}
\begin{teile}
\teil Stelle die Gleichung für den Zeitpunkt auf, zu dem der Ball am Boden ist.
\teil Löse die Gleichung durch Ausklammern.
\teil Deute beide Lösungen im Sachzusammenhang.
\end{teile}
\end{aufgabe}
